
\subsection{5.1. The representations carry batch information}\label{the-representations-carry-batch-information}

In GSE174188 the three representations classified disease at AUC 0.982, 0.977 and 0.975. The same representations identified processing wave, one against the rest, at 0.9997, 0.999 and 0.997. In GSE135779 disease AUCs were 0.870, 0.933 and 0.945, while the best single-batch AUCs were 0.993, 0.961 and 0.961 (Figure 5A).

This shows that the representations contain batch information. It shows nothing more. A representation that encodes batch can still carry disease information, and a batch-aware representation need not produce a batch-driven decision rule. Batch predictability is therefore not a validity check on its own, and we report it below only with that scope attached.

\subsection{5.2. Containing batch information is not the same as using it}\label{containing-batch-information-is-not-the-same-as-using-it}

GSE135779 provides the counterexample directly. The best one-against-rest batch AUC was 0.993 for Geneformer and 0.961 for both pseudobulk representations, which is close to perfect batch recognition. Yet batch alone predicted disease at AUC 0.499, and fold-contained batch residualisation changed disease AUC by +0.002, −0.002 and −0.009 (Figure 5B, 6C). The information was there, and removing it changed nothing measurable at this sample size - the three shifts are well inside the split-to-split spread, so we can say the classifier's dependence on batch is small, not that it is zero. That is enough to break the inference. Going from ``the representation knows the batch'' to ``the classifier uses the batch'' does not hold on these data.

\subsection{5.3. Residualisation and restriction answer different questions}\label{residualisation-and-restriction-answer-different-questions}

Restriction asks whether the separation survives among donors collected under overlapping conditions. Residualisation deletes the part of the data that a chosen set of design variables can predict. When design and diagnosis are collinear these are not the same question, and we compare them to show they are not interchangeable, not to rank them.

In GSE174188, ridge residualisation on processing-wave fractions reduced the three representations from 0.982, 0.977 and 0.975 to 0.714, 0.747 and 0.735. Restriction to the one large near-balanced wave, compared against size- and label-matched random donor subsets, cost only 0.014 to 0.049 (Figure 6B). The paired difference between the two losses was 0.186 to 0.255 across six comparisons, with every descriptive interval above zero.

In GSE135779, removing the single all-case batch left 36 donors with both labels in every remaining batch and changed matched performance by at most 0.012. Five ways of removing the recorded variables give five different answers for the same cohort (Figure 6A). Projecting out the complete set with ridge collapsed all three representations to 0.516, 0.522 and 0.520. A random-forest residualiser removed less (down to 0.748, 0.787 and 0.795), and overlap weighting and batch location adjustment less still (retaining 0.859 to 0.937). Every full-set removal took away far more than the batch-only control.

The algebra explains this without needing a technical story. When design predicts the diagnosis, the part of the data that design can predict necessarily includes diagnosis-related variation. Removing it removes disease signal regardless of where that signal came from. A large loss after residualisation is therefore \textbf{not} evidence that the original AUC was spurious. It is equally consistent with removing a shortcut and with removing real biology, and the size of the loss does not distinguish them.

A negative control settles part of what residualisation loss actually measures. In the CMV cohort no association is detectable between the full recorded set and the diagnosis: design-only AUC 0.483 under this pipeline, p(V\_D) 0.284 to 0.353 across five seeds. Residualising its representations on those variables still costs 0.282 AUC, from 0.845 down to 0.563 (Figure 6C). There is nothing there to remove, and removing it costs a quarter of the separation.

Our first explanation was arithmetic: 89 one-hot columns for 108 donors, so the projection takes out most of any sample. We tested that rather than asserting it. We ran a simulation in which the design is independent of the diagnosis by construction, with the unadjusted AUC matched to this cohort. A balanced design matrix of the same width (92 columns) removes only 0.045. The loss is also not monotone in width: it peaks at 0.229 for 64 columns and falls away on either side. One built to the same ragged shape as the real batch-by-pool crossing - many levels holding one or two donors - removes 0.091, twice what balanced blocks of similar width remove (Figure S2). Width and raggedness account for part of the 0.282 and not for all of it. So the claim we can support is narrower than the one we first wrote: residualisation loss depends on properties of the design matrix that have nothing to do with confounding, and its size therefore cannot be read as a measure of technical contamination. We cannot fully account for this particular cohort's loss, and we report the gap rather than close it with a story.

It is tempting to read the pair of values from residualisation and restriction as a range bounding the true biological effect. It is not one. There is no proof that the two estimates bracket any underlying quantity. They are two estimates made under two different sets of assumptions, and their disagreement is informative precisely because neither is identified.

\subsection{5.4. Cell-type composition shows the same pattern}\label{cell-type-composition-shows-the-same-pattern}

Aggregating the same GSE174188 CD4 cells into naive, effector-memory, regulatory and unassigned fractions gave disease AUC 0.899 to 0.915. The same four-number summaries predicted processing wave at 0.803 to 0.872. Wave residualisation reduced them to 0.689 to 0.709 (Figure S5A), and within-wave restriction against matched subsets cost 0.027 to 0.062, with every interval crossing zero (Figure S5B). A four-dimensional biological summary therefore reproduces the whole pattern. That places the problem in the cohort rather than in foundation-model embeddings, and is consistent with composition being a strong patient stratification baseline in its own right {[}26{]}.

\subsection{5.5. Transfer to an independent cohort does not fix the problem}\label{transfer-to-an-independent-cohort-does-not-fix-the-problem}

Restricting the training set to one processing wave is the obvious remedy: train where the design is uniform, then test elsewhere. It does not work. Training on all 261 GSE174188 donors and testing in the independent 26-donor GSE285773 cohort gave 0.900, 0.944 and 0.969. Restricting the source to the dominant wave lowered every one of them, and did no better than a random subset of the same size (Figure 7A). Within GSE174188, training on waves 2 and 3 and testing on wave 4 gave 0.842, 0.811 and 0.806, and the reverse direction 0.919, 0.926 and 0.821. More varied source donors transferred better than the apparently cleaner restricted source.

Transfer in the other direction, from the 26-donor cohort to the 261-donor target under strict source-only rules, gave 0.884 {[}0.843-0.920{]}, 0.919 {[}0.884-0.946{]} and 0.926 {[}0.895-0.953{]}. Both pseudobulk representations beat Geneformer after paired DeLong testing {[}23{]} and Benjamini-Hochberg adjustment {[}24{]} (q = 0.0151, q = 0.000603).

Whether a more expressive donor-level aggregator changes this is a separate question from the one this paper asks, and we report it in the supplement rather than the main text. Briefly: three learned pooling methods on the same frozen embeddings never beat ordinary mean pooling, and every significant paired difference was negative (Figure S6, Table S5). Adding capacity fitted the source without producing a representation that travels.

\section{6. Five checks, and what each one can show}\label{five-checks-and-what-each-one-can-show}

These are not a validity standard. We have not shown that they are sufficient, that they are minimal, or that passing them licenses a causal claim. They are five checks that look at different arrows in Figure 1 and that fail in different ways. Their value is that they can disagree with each other.

\begin{longtable}[]{@{}
  >{\raggedright\arraybackslash}p{(\columnwidth - 6\tabcolsep) * \real{0.2500}}
  >{\raggedright\arraybackslash}p{(\columnwidth - 6\tabcolsep) * \real{0.2500}}
  >{\raggedright\arraybackslash}p{(\columnwidth - 6\tabcolsep) * \real{0.2500}}
  >{\raggedright\arraybackslash}p{(\columnwidth - 6\tabcolsep) * \real{0.2500}}@{}}
\toprule\noalign{}
\begin{minipage}[b]{\linewidth}\raggedright
\#
\end{minipage} & \begin{minipage}[b]{\linewidth}\raggedright
Check
\end{minipage} & \begin{minipage}[b]{\linewidth}\raggedright
Shows
\end{minipage} & \begin{minipage}[b]{\linewidth}\raggedright
Does \textbf{not} show
\end{minipage} \\
\midrule\noalign{}
\endhead
\bottomrule\noalign{}
\endlastfoot
1 & Design-only AUC by variable group, cross-fitted V\_D, p(V\_D) against the cohort's own null, overlap counts & Whether recorded design predicts the diagnosis, and which variables do & That the classifier uses design; anything about unrecorded structure \\
2 & Complete-pipeline permutation, both free \textbf{and} collection-preserving & Free: pipeline validity, no leakage. Collection-preserving: that accuracy beats what the collection strata alone can produce & Collection-preserving conditions on collection variables only, so it does not cover sample-quality or demographic association inside a stratum (Section 3.4); neither is defined when no stratum holds both labels; neither identifies a mechanism \\
3 & How well the representation predicts design & That the representation contains collection information & That the diagnosis decision rule uses it --- Section 5.2 is a counterexample \\
4 & Residualisation \textbf{paired with} restriction and matched donor controls & How far two non-equivalent adjustments disagree & Which adjustment is right; bounds on a biological effect \\
5 & Source-only external target & Whether patient ranking survives a change of collection process & Calibration or clinical usefulness; and it carries much less weight when the target was collected the same way, since a shared bias reproduces rather than cancels (Section 3.3) \\
\end{longtable}

Checks 1 and 2 cost minutes of CPU on a donor-level table and are the ones we would run first on any new cohort. Check 5 is the only one that speaks to transport, and no combination of the internal checks substitutes for it.

\textbf{What a clinician should take from this.} A published patient-level classifier is worth taking seriously as a candidate diagnostic when three things are reported together. The first is the design-only AUC for the development cohort. The second is both permutation tests, or a statement that the collection-preserving one is undefined. The third is an external result on a cohort collected through a different route. If the paper reports only a cross-validated AUC, the number is uninterpretable - not wrong, but not yet a claim about disease. If it reports an external validation on a cohort assembled the same way as the first, that still demonstrates reproducibility, but it does not rule out a bias shared by both.

Applied to our nine comparisons, the checks give three groups rather than a pass and a fail. In the first, no association with the recorded metadata is detectable, and internal accuracy is not undermined by anything we can see - which is weaker than saying it is biological, since unrecorded structure remains possible; only CMV is here. In the second, the metadata predicts the diagnosis strongly and yet the collection-preserving permutation is still beaten, so the accuracy exceeds what the collection strata alone explain; both lupus comparisons and all three COVID-19 cohorts are here. In the third, the collection-preserving permutation is either not beaten or not defined, and no internal analysis can settle the question; sepsis-versus-COVID-19 and influenza are here. Only the third group is a stop condition, and it is a statement about the cohort, not about the method being evaluated.

\subsection{6.1. Three cohorts taken through the tree step by step}\label{three-cohorts-taken-through-the-tree-step-by-step}

Figure 8 traces four comparisons through the decision tree with the evidence for each branch stated. Three are given in full here; the fourth, sepsis versus COVID-19, is described in Section 4.4. The strata and the design matrix are both imported from the screen rather than rebuilt for the walkthrough, so Q2 and Q4 ask about the same partition and the same columns.

\textbf{COMBAT influenza (n = 21) stops at Q0.} Its stratum variable is the recording institute, which gives two strata, \emph{neither of which contains both a case and a control} (Figure S4). The collection-preserving permutation therefore has nothing to shuffle and is undefined. The comparison therefore cannot be answered with these data and no route in the tree is available. The frozen pipeline also scores far below the tuned one here - 0.27 to 0.58 against 0.98 to 1.00 across seeds. That gap is reported but is not a gate anywhere in this paper, and it is not what stops the trace. The stopping condition is the absence of a stratum holding both labels. Note what the free permutation test alone would have said: p = 0.299 to 0.966 across the five seeds, which reads as ``no signal'' and is the wrong description.

\textbf{CMV HIHA (n = 108) reaches Route A.} The collection-preserving permutation is defined - 46 strata from batch and pool, 21 of them containing both labels - and both permutation tests reject at the 1/1001 floor. At Q3 the association with the recorded metadata is weak (design-only AUC 0.483 under the walkthrough's pipeline, 0.470 to 0.501 in the main screen, not significant against its own permutation null in any seed), so the comparison takes the adjustment route without reaching the restriction questions. It is the only one of our nine comparisons that does.

\textbf{COVID-19 Ren (n = 182) reaches Route B2.} Both permutation tests are beaten. The association with the recorded metadata is strong: design-only AUC 0.840 under the walkthrough's fixed-penalty pipeline and 0.805 to 0.832 under the main one, both far above chance. Mixed strata exist, so Q4 passes and the comparison reaches the residualisation-versus-restriction question. Residualising the 26-column design matrix costs 0.100, from 0.948 down to 0.849 AUC. Restriction gives the opposite sign: within the one stratum large enough to cross-validate, AUC is 0.972 against 0.749 in size- and composition-matched random subsets. The two adjustments disagree in direction. That is Route B2, and under the revised tree the trace stops there, with two estimates reported and no bound claimed.

\textbf{The walkthrough also exposes a weakness in the tree itself.} At the stratum definition the collection-preserving permutation uses, restriction is barely available in these cohorts. CMV has 21 label-mixed strata but \emph{none} with five donors in each class. COVID-19 Ren has exactly one usable stratum, of 18 donors, whose matched-subset comparison carries a standard deviation of 0.117. Route B is drawn as though restriction were a practical alternative to residualisation. On real cohorts with realistically fine strata it frequently is not, and the honest output in those cases is Route C - redesign - rather than a restricted estimate computed on eighteen donors. We state this rather than loosening the stratum definition until Route B becomes available, which would make the tree unfalsifiable.

\section{Discussion}\label{discussion}

\section{7. What this means for reading a published classifier}\label{what-this-means-for-reading-a-published-classifier}

Design-diagnosis association in public single-cell cohorts is measurable, graded, and traceable to specific variables. Across nine comparisons from five widely reused cohorts it runs from a design-only AUC of 0.470 to 1.000, which is more than half the usable scale, and a cohort's position is stable across random seeds. The \emph{raw} residual label variance, by contrast, is not comparable across cohorts without its own null. That matters in practice: an in-sample R² ranks the least affected cohort in our set as the third most affected, purely because its design matrix is wide relative to its donor count.

Grouping the design variables matters more than the single number. Two COVID-19 cohorts whose design-only AUCs differ by 0.02 are affected through different variables, one through library quality and one through recruiting hospital. Restriction to balanced strata is available for one and not the other. A quality-driven association may be partly correctable by harmonised preprocessing; a hospital-driven one is not. Reporting one number hides the only part of the answer a reader can act on.

The disagreement between the two permutation tests is, in our view, the finding with the widest reach beyond this dataset. Complete-pipeline label permutation is standard practice and is a good test of what it tests. But it cannot test whether a classifier exploits the collection process, because it destroys that process as part of the shuffle. On the sepsis-versus-COVID-19 comparison the two tests reach opposite conclusions, and the conventional one is the permissive one. Any study that uses label permutation to argue that a patient-level classifier is valid should report the collection-preserving version alongside it, or state that no mixed stratum exists and that version is undefined. That statement is itself the more informative outcome.

The replacement test has a failure mode of its own, and we would rather publish the number than the caveat. It conditions on collection strata, not on the whole recorded matrix, so a sample-quality variable that stays associated with the diagnosis inside a stratum can carry a classifier past it with no disease effect present. Section 3.4 measures how often. A significant collection-preserving result therefore rules out one specific explanation - these strata - and not the general one.

The simulation adds a caution about external validation that we did not anticipate. When the external cohort shares the source's collection mechanism, external AUC rises with confounding rather than falling. An external result is evidence about transport only in proportion to how independent the second collection process really is, and being a different accession number does not establish independence.
